% ============================================================
%  KLINISK DANSK — bilingual reference/glossary (PDF)
%  Kapitel 1: Kredsløb & blodtryk
%  Kompilér: pdflatex main.tex  (kør 2x for indholdsfortegnelse)
%  NB: babel er bevidst udeladt (txtbabel.def-fejl på denne MiKTeX).
% ============================================================
\documentclass[11pt,a4paper]{report}
\usepackage[utf8]{inputenc}
\usepackage[T1]{fontenc}
\usepackage{lmodern}
\usepackage{textcomp}
\usepackage{amssymb}
\usepackage{pifont}
\usepackage[margin=2.4cm]{geometry}
\usepackage{longtable}
\usepackage{booktabs}
\usepackage{array}
\usepackage{xcolor}
\usepackage{tcolorbox}
\tcbuselibrary{skins,breakable}
\usepackage{titlesec}
\usepackage{enumitem}
\usepackage{fancyhdr}
\usepackage[hidelinks]{hyperref}

% ---- farver (matcher web-appen) ----
\definecolor{brand}{HTML}{0E7C86}
\definecolor{brandd}{HTML}{0B5F67}
\definecolor{en}{HTML}{7A5CC0}
\definecolor{good}{HTML}{0F8A4D}
\definecolor{goodbg}{HTML}{E6F6EE}
\definecolor{bad}{HTML}{C0392B}
\definecolor{badbg}{HTML}{FBECEA}
\definecolor{ink}{HTML}{12222E}
\definecolor{muted}{HTML}{5A6B78}
\definecolor{line}{HTML}{D7E0E8}

\color{ink}

% ---- overskrifter ----
\titleformat{\chapter}[hang]{\color{brandd}\bfseries\Huge}{\thechapter\quad}{0pt}{}
\titlespacing*{\chapter}{0pt}{0pt}{18pt}
\titleformat{\section}{\color{brandd}\bfseries\Large}{\thesection}{0.6em}{}
\titleformat{\subsection}{\color{brand}\bfseries\large}{\thesubsection}{0.6em}{}

% ---- sidehoved/-fod ----
\pagestyle{fancy}\fancyhf{}
\renewcommand{\headrulewidth}{0.4pt}
\fancyhead[L]{\small\color{muted}Klinisk Dansk}
\fancyhead[R]{\small\color{muted}\nouppercase{\leftmark}}
\fancyfoot[C]{\small\color{muted}\thepage}

% ---- hjælpe-makroer ----
\newcommand{\eng}[1]{\textcolor{en}{\itshape #1}}       % engelsk gloss
\newcommand{\da}[1]{\textbf{#1}}
% glossar-tabel: DA | ordklasse | EN | eksempel
\newenvironment{glos}{%
  \small\renewcommand{\arraystretch}{1.35}
  \begin{longtable}{@{}>{\raggedright\arraybackslash}p{3.3cm}
                        >{\raggedright\arraybackslash}p{1.7cm}
                        >{\raggedright\arraybackslash}p{3.2cm}
                        >{\raggedright\arraybackslash}p{6.0cm}@{}}
  \toprule
  \textbf{Dansk} & \textbf{Ordklasse} & \eng{English} & \textbf{Eksempel} \\ \midrule
  \endhead
}{\bottomrule\end{longtable}}

% god/dårlig-boks
\newtcolorbox{daarligbox}{colback=badbg,colframe=bad,boxrule=0.4pt,left=8pt,right=8pt,top=6pt,bottom=6pt,
  breakable,fonttitle=\bfseries\color{bad},title={\ding{55}\ Dårligt}}
\newtcolorbox{godbox}{colback=goodbg,colframe=good,boxrule=0.4pt,left=8pt,right=8pt,top=6pt,bottom=6pt,
  breakable,fonttitle=\bfseries\color{good},title={\checkmark\ Godt}}
\newtcolorbox{hvorforbox}{colback=white,colframe=line,boxrule=0.4pt,left=8pt,right=8pt,top=6pt,bottom=6pt,
  breakable,fonttitle=\bfseries\color{brandd},title={Hvorfor}}
\newtcolorbox{notebox}{colback=white,colframe=brand,boxrule=0.6pt,left=8pt,right=8pt,top=6pt,bottom=6pt,breakable}

\begin{document}

% ---------------- forside ----------------
\begin{titlepage}
\centering
\vspace*{3cm}
{\color{brandd}\Huge\bfseries Klinisk Dansk\par}
\vspace{0.4cm}
{\LARGE Sygeplejesprog — bilingual referenceværk\par}
\vspace{0.2cm}
{\large\color{muted}Danish for nursing / healthcare · niveau B2–C1\par}
\vspace{2.5cm}
{\color{brand}\rule{6cm}{2pt}\par}
\vspace{0.6cm}
{\Large\bfseries Fire dage --- kredsløb, respiration, smerte \& medicin\par}
\vspace{0.3cm}
{\large\color{muted}Dag 1 Kredsløb \& blodtryk · Dag 2 Respiration \& saturation\\
Dag 3 Smerte \& neuro · Dag 4 Medicin \& ordination\par}
\vfill
{\color{muted}Ordforråd · dokumentationssprog · patientnotater · øvelser med facit\par}
\end{titlepage}

\tableofcontents
\newpage

% ================================================================
\chapter{Kredsløb \& blodtryk}

Denne dag dækker hjertet, karrene og blodtrykket --- og hvordan du beskriver det, du ser
og måler, i \textbf{objektivt fagsprog}. Bogen er tosproget: dansk er det primære sprog,
og hver term og hvert eksempel har en engelsk oversættelse (\eng{English in violet italics}).

% ---------------------------------------------------------------
\section{Ordforråd}

\subsection{Substantiver --- organer, symptomer, måleværdier}
\begin{glos}
\da{hjertet} & substantiv & \eng{the heart} & Hjertet pumper iltet blod ud i det store kredsløb. \\
\da{blodkar / kar} & subst. & \eng{blood vessel} & Karrene i benene er svære at palpere pga. ødem. \\
\da{arterie (pulsåre)} & subst. & \eng{artery} & Blodet i en arterie er iltrigt og lyserødt. \\
\da{vene (blodåre)} & subst. & \eng{vein} & Der anlægges et perifert venekateter (PVK) i venstre underarm. \\
\da{kapillær (hårkar)} & subst. & \eng{capillary} & Kapillærresponsen (CR) er under 2 sekunder --- normal. \\
\da{hjerteklap} & subst. & \eng{heart valve} & Der høres en mislyd fra en af hjerteklapperne. \\
\da{blodtryk (BT)} & subst. & \eng{blood pressure} & BT måles til 138/86 mmHg i højre arm. \\
\da{det systoliske tryk} & subst. & \eng{systolic pressure} & Det systoliske tryk er forhøjet. \\
\da{det diastoliske tryk} & subst. & \eng{diastolic pressure} & Det diastoliske tryk ligger stabilt omkring 80. \\
\da{puls} & subst. & \eng{pulse} & Pulsen er 102 og regelmæssig. \\
\da{hjertefrekvens (HF)} & subst. & \eng{heart rate} & HF falder til 68 efter hvile. \\
\da{iltmætning / saturation (SAT)} & subst. & \eng{oxygen saturation} & Saturation 96\,\% på atmosfærisk luft. \\
\da{brystsmerter} & subst. & \eng{chest pain} & Pt. angiver trykkende brystsmerter med udstråling til venstre arm. \\
\da{hjertebanken / palpitationer} & subst. & \eng{palpitations} & Pt. beskriver anfald af hjertebanken i hvile. \\
\da{ødem} & subst. & \eng{oedema / swelling} & Der ses bilateralt ankelødem, gradtrykbart. \\
\da{blodprop (trombe)} & subst. & \eng{blood clot / thrombus} & Mistanke om blodprop i venstre underben (DVT). \\
\da{hjerteinfarkt (AMI)} & subst. & \eng{myocardial infarction} & Symptomerne rejser mistanke om akut hjerteinfarkt. \\
\da{ordination} & subst. & \eng{order / prescription} & Nitroglycerin gives efter ordination. \\
\da{kredsløbet} & subst. & \eng{the circulation} & Pt. er kredsløbsmæssigt stabil. \\
\end{glos}

\subsection{Verber --- sygepleje-handlinger}
\begin{glos}
\da{at observere} & verbum & \eng{to observe} & Pt. observeres for ændringer i bevidsthedsniveau. \\
\da{at måle} & verbum & \eng{to measure} & Der måles BT, puls og saturation hver time. \\
\da{at palpere} & verbum & \eng{to palpate} & Radialispulsen palperes kraftig og regelmæssig. \\
\da{at auskultere} & verbum & \eng{to auscultate} & Hjertet auskulteres uden mislyde. \\
\da{at monitorere} & verbum & \eng{to monitor} & Pt. monitoreres kontinuerligt med EKG. \\
\da{at administrere} & verbum & \eng{to administer} & Furosemid 40 mg administreres i.v. efter ordination. \\
\da{at ordinere} & verbum & \eng{to prescribe / order} & Lægen ordinerer ilt 2 l/min på næsekateter. \\
\da{at dokumentere} & verbum & \eng{to document} & Observationerne dokumenteres i sygeplejejournalen. \\
\da{at registrere} & verbum & \eng{to record} & Værdierne registreres i observationsskemaet. \\
\da{at angive} & verbum & \eng{to report / state} & Pt. angiver smerter svarende til 6 på VAS-skalaen. \\
\da{at klage over} & verbum & \eng{to complain of} & Pt. klager over svimmelhed ved sengekanten. \\
\da{at desaturere} & verbum & \eng{to desaturate} & Pt. desaturerer til 88\,\% ved gang på gangen. \\
\end{glos}

\subsection{Adjektiver --- beskrivelser af symptomer og værdier}
\begin{glos}
\da{regelmæssig / uregelmæssig} & adjektiv & \eng{regular / irregular} & Pulsen er uregelmæssig --- mistanke om atrieflimren. \\
\da{forhøjet} & adjektiv & \eng{elevated} & Blodtrykket er forhøjet gennem hele vagten. \\
\da{hypertensiv} & adjektiv & \eng{hypertensive} & Pt. er hypertensiv med BT 178/104. \\
\da{hypotensiv} & adjektiv & \eng{hypotensive} & Pt. bliver hypotensiv ved skift til stående (ortostatisk). \\
\da{takykard} & adjektiv & \eng{tachycardic} & Pt. er takykard med HF 118 i hvile. \\
\da{bradykard} & adjektiv & \eng{bradycardic} & Pt. er bradykard med HF 46 under søvn. \\
\da{perifer} & adjektiv & \eng{peripheral} & Den perifere cirkulation i tæerne er nedsat. \\
\da{kold og klam} & adjektiv & \eng{cold and clammy} & Huden er kold og klam, pt. er bleg. \\
\da{marmoreret} & adjektiv & \eng{mottled} & Huden på knæene er marmoreret --- tegn på dårlig cirkulation. \\
\da{jagende / skarp} & adjektiv & \eng{sharp / stabbing} & Pt. beskriver en jagende smerte ved indånding. \\
\da{dump} & adjektiv & \eng{dull} & En dump, murrende smerte bag brystbenet. \\
\da{murrende} & adjektiv & \eng{aching / nagging} & Pt. angiver murrende ubehag i brystet. \\
\da{trykkende} & adjektiv & \eng{pressing / squeezing} & Trykkende brystsmerter, som ved en snørende fornemmelse. \\
\da{bleg} & adjektiv & \eng{pale} & Pt. er bleg og forpint ved ankomst. \\
\da{cyanotisk} & adjektiv & \eng{cyanotic} & Læberne er let cyanotiske. \\
\end{glos}

% ---------------------------------------------------------------
\section{Dokumentationssprog}

\subsection{Grundregler}
\begin{enumerate}[leftmargin=1.4em,itemsep=4pt]
\item \textbf{Skriv objektivt} --- beskriv hvad du ser, hører og måler, ikke hvad du tror eller føler.\\ \eng{Write objectively --- what you see, hear and measure, not what you think or feel.}
\item \textbf{Ingen personlig vurdering.} Undgå ord som ``virker'', ``ser ud til'', ``er sød/besværlig''.\\ \eng{No personal judgement. Avoid ``seems'', ``looks like'', ``is nice/difficult''.}
\item \textbf{Adskil hvad patienten SIGER fra hvad du OBSERVERER.} Brug ``pt.\ angiver \ldots'' for det subjektive.\\ \eng{Separate what the patient SAYS from what you OBSERVE. Use ``pt.\ reports \ldots''.}
\item \textbf{Brug tal og enheder:} BT 138/86 mmHg, puls 102, SAT 96\,\%, VAS 6/10.\\ \eng{Use numbers and units.}
\item \textbf{Angiv altid betingelsen} for en måling: ``SAT 92\,\% på atmosfærisk luft'' vs.\ ``på 2 l ilt''.\\ \eng{Always state the condition of a measurement (room air vs.\ oxygen).}
\item \textbf{Skriv kort og telegramagtigt.} Forkortelser: pt.\ (patient), i.v., p.n.\ (efter behov).\\ \eng{Write short and telegraphic; standard abbreviations.}
\end{enumerate}

\subsection{Dårlig vs.\ god formulering}

\begin{daarligbox}Pt.\ virker dårlig og har vist lidt ondt i brystet.\end{daarligbox}
\begin{godbox}Pt.\ angiver trykkende brystsmerter VAS 6/10 med udstråling til venstre arm. BT 165/98, puls 102 regelmæssig, SAT 95\,\% på atm.\ luft, hud kold og klam.\end{godbox}
\begin{hvorforbox}``Virker dårlig'' og ``vist lidt'' er vurdering og gætteri. Den gode note adskiller det subjektive (pt.\ angiver) fra objektive målte værdier med tal og enheder.\end{hvorforbox}
\medskip

\begin{daarligbox}Gav ham hans hjertemedicin.\end{daarligbox}
\begin{godbox}Metoprolol 50 mg administreret p.o.\ kl.\ 08.00 efter ordination.\end{godbox}
\begin{hvorforbox}Dokumentér præparat, dosis, administrationsvej, tidspunkt og at det skete efter ordination --- aldrig bare ``medicinen''.\end{hvorforbox}
\medskip

\begin{daarligbox}Pt.\ er sur og vil ikke samarbejde om blodtrykket.\end{daarligbox}
\begin{godbox}Pt.\ afviser BT-måling kl.\ 10 og angiver, at han ``vil have ro''. BT ikke målt. Informeret om formålet, forsøges igen kl.\ 12.\end{godbox}
\begin{hvorforbox}``Sur'' og ``vil ikke samarbejde'' er en negativ vurdering. Beskriv i stedet handlingen (afviser), patientens egne ord i citat, og hvad du gør ved det.\end{hvorforbox}
\medskip

\begin{daarligbox}Saturationen var fin.\end{daarligbox}
\begin{godbox}SAT 97\,\% på atmosfærisk luft, upåvirket respiration, RF 16.\end{godbox}
\begin{hvorforbox}``Fin'' siger intet. Angiv den målte værdi, betingelsen (atm.\ luft/ilt) og gerne respirationsfrekvensen.\end{hvorforbox}
\medskip

\begin{daarligbox}Benene er lidt hævede, tror det er væske.\end{daarligbox}
\begin{godbox}Bilateralt ankelødem, gradtrykbart. Vægt +1,3 kg ift.\ i går. Lægen orienteret.\end{godbox}
\begin{hvorforbox}Erstat ``tror'' og ``lidt'' med en objektiv beskrivelse (bilateralt, gradtrykbart), et måltal (vægtændring) og handlingen (lægen orienteret).\end{hvorforbox}

% ---------------------------------------------------------------
\section{Patientnotat}

\subsection{Eksempel --- indlæggelse, medicinsk afdeling}
\begin{notebox}
Pt.\ 74-årig mand indlagt kl.\ 21.40 med trykkende brystsmerter opstået i hvile, VAS 7/10,
med udstråling til venstre skulder. Pt.\ bleg og klamtsvedende. BT 178/102 mmHg, puls 108
uregelmæssig, SAT 93\,\% på atmosfærisk luft, RF 22, temp.\ 37,1\,\textdegree C. Perifert kølige
ekstremiteter, kapillærrespons 3 sek. Nitroglycerin 0,5 mg gives sublingualt efter ordination
kl.\ 21.55; smerte falder til VAS 3/10. Ilt 2 l/min ordineret og påbegyndt; SAT stiger til
96\,\%. Pt.\ monitoreres med EKG. Lægen tilser.

\medskip
{\color{muted}\small\eng{74-year-old man admitted 21:40 with pressing chest pain arising at rest, VAS 7/10,
radiating to the left shoulder. Pt.\ pale and clammy. BP 178/102 mmHg, pulse 108 irregular,
SAT 93\% on room air, RR 22, temp 37.1\,\textdegree C. Peripherally cool extremities, capillary
refill 3 sec. Nitroglycerin 0.5 mg given sublingually as prescribed at 21:55; pain drops to
VAS 3/10. Oxygen 2 l/min prescribed and started; SAT rises to 96\%. Pt.\ monitored with ECG.
Doctor attends.}}
\end{notebox}

\subsection{Øvelser}
\begin{enumerate}[leftmargin=1.5em,itemsep=8pt]
\item \textbf{Omskriv til objektivt fagsprog:} ``Den gamle mand så virkelig dårlig ud og havde det skidt med brystet.''
\item \textbf{Omskriv:} ``Gav ham noget under tungen og det hjalp på smerten.''
\item \textbf{Udfyld verbet:} ``SAT 93\,\% på atmosfærisk luft. Ilt 2 l/min \underline{\hspace{2.5cm}} og påbegyndt; SAT stiger til 96\,\%.''
\item \textbf{Udfyld verbet:} ``Metoprolol 50 mg \underline{\hspace{2.5cm}} p.o.\ efter ordination kl.\ 08.00.''
\item \textbf{Find vurderingen}, der ikke hører hjemme: ``Pt.\ angiver kvalme. BT 138/84. Pt.\ er en smule krævende. SAT 97\,\%.''
\end{enumerate}

\subsection{Facit}
\begin{enumerate}[leftmargin=1.5em,itemsep=5pt]
\item Pt.\ er bleg og forpint, angiver trykkende brystsmerter VAS 7/10 med udstråling til venstre skulder.
\item Nitroglycerin 0,5 mg administreret sublingualt efter ordination; smerte falder fra VAS 7/10 til 3/10.
\item \textbf{ordineret} (lægen ordinerer ilt).
\item \textbf{administreret} (sygeplejersken administrerer medicinen).
\item ``Pt.\ er en smule krævende'' --- en personlig, negativ vurdering uden observation bag.
\end{enumerate}

% ================================================================
\chapter{Respiration \& saturation}

Denne dag handler om vejrtrækningen: respirationsfrekvens, iltmætning, iltbehandling og de
lyde og symptomer, du observerer --- dyspnø, hoste, ekspektorat, rallen og hvæsen. Kernen er
altid at angive \textbf{iltbetingelsen} sammen med saturationen.

\section{Ordforråd}
\subsection{Substantiver}
\begin{glos}
\da{respiration / vejrtrækning} & subst. & \eng{respiration / breathing} & Vejrtrækningen er overfladisk og anstrengt. \\
\da{respirationsfrekvens (RF)} & subst. & \eng{respiratory rate} & RF 24 i hvile --- let forhøjet. \\
\da{åndenød / dyspnø} & subst. & \eng{dyspnoea} & Pt.\ angiver tiltagende åndenød ved gang. \\
\da{iltmætning / saturation (SAT)} & subst. & \eng{oxygen saturation} & SAT 91\,\% på atm.\ luft, stiger til 96\,\% på 2 l ilt. \\
\da{luftvej} & subst. & \eng{airway} & Pt.\ har fri luftvej og taler i hele sætninger. \\
\da{ekspektorat / opspyt} & subst. & \eng{sputum} & Gulligt, purulent ekspektorat. \\
\da{hoste} & subst. & \eng{cough} & Produktiv hoste med sejt slim. \\
\da{ilt / oxygen (O$_2$)} & subst. & \eng{oxygen} & Ilt 2 l/min gives på næsekateter efter ordination. \\
\da{næsekateter} & subst. & \eng{nasal cannula} & Ilten leveres via næsekateter. \\
\da{iltmaske} & subst. & \eng{oxygen mask} & Skiftet til iltmaske med reservoir, 10 l/min. \\
\da{lungeødem} & subst. & \eng{pulmonary oedema} & Rallen ved basis giver mistanke om lungeødem. \\
\da{lungebetændelse (pneumoni)} & subst. & \eng{pneumonia} & Pt.\ behandles for pneumoni med i.v.\ antibiotika. \\
\da{KOL} & subst. & \eng{COPD} & Pt.\ er kendt med svær KOL. \\
\da{cyanose} & subst. & \eng{cyanosis} & Let central cyanose omkring læberne. \\
\end{glos}

\subsection{Verber}
\begin{glos}
\da{at trække vejret} & verbum & \eng{to breathe} & Pt.\ trækker vejret anstrengt med brug af hjælpemuskler. \\
\da{at indånde / inhalere} & verbum & \eng{to inhale} & Medicinen inhaleres via forstøverapparat. \\
\da{at udånde / ekspirere} & verbum & \eng{to exhale} & Der ses forlænget ekspiration, typisk for KOL. \\
\da{at hoste} & verbum & \eng{to cough} & Pt.\ hoster produktivt gennem natten. \\
\da{at ekspektorere} & verbum & \eng{to expectorate} & Pt.\ ekspektorerer sejt, gulligt slim. \\
\da{at auskultere} & verbum & \eng{to auscultate} & Lungerne auskulteres med rallen ved begge baser. \\
\da{at desaturere} & verbum & \eng{to desaturate} & Pt.\ desaturerer til 87\,\% ved gang til toilettet. \\
\da{at iltbehandle} & verbum & \eng{to give oxygen therapy} & Pt.\ iltbehandles med 2 l/min efter ordination. \\
\da{at hyperventilere} & verbum & \eng{to hyperventilate} & Pt.\ hyperventilerer i forbindelse med angst. \\
\end{glos}

\subsection{Adjektiver}
\begin{glos}
\da{overfladisk} & adjektiv & \eng{shallow} & Vejrtrækningen er overfladisk og hurtig. \\
\da{anstrengt / besværet} & adjektiv & \eng{laboured} & Anstrengt respiration med indtrækninger. \\
\da{rallende} & adjektiv & \eng{crackling / rattling} & Rallende vejrtrækning, hørbar uden stetoskop. \\
\da{pibende / hvæsende} & adjektiv & \eng{wheezing} & Pibende, hvæsende ekspiration ved auskultation. \\
\da{produktiv / tør} & adjektiv & \eng{productive / dry} & Hosten er nu produktiv, tidligere tør. \\
\da{purulent / klar} & adjektiv & \eng{purulent / clear} & Ekspektoratet er purulent og gulligt. \\
\da{takypnø} & adjektiv & \eng{tachypnoeic} & Pt.\ er takypnø med RF 28 i hvile. \\
\da{cyanotisk} & adjektiv & \eng{cyanotic} & Læber og negle er let cyanotiske. \\
\end{glos}

\section{Dokumentationssprog}
\subsection{Grundregler}
\begin{enumerate}[leftmargin=1.4em,itemsep=4pt]
\item \textbf{Angiv ALTID iltbetingelsen} ved en saturation: ``SAT 92\,\% på atmosfærisk luft'' vs.\ ``på 2 l ilt''.
\item \textbf{Beskriv respirationen} med RF + karakter (overfladisk/anstrengt/upåvirket), ikke ``trækker vejret fint''.
\item \textbf{Karakterisér hoste} (produktiv/tør) og ekspektorat (mængde, farve, konsistens).
\item \textbf{Skriv objektive lyde:} ``rallen ved basis'', ``pibende ekspiration'' --- ikke ``lyder ikke godt''.
\item \textbf{Dokumentér effekten} af iltbehandling: ``SAT stiger fra 89\,\% til 95\,\% på 2 l ilt''.
\item \textbf{Ved KOL:} følg det ordinerede SAT-mål (ofte 88--92\,\%); for høj ilt kan skade.
\end{enumerate}

\subsection{Dårlig vs.\ god formulering}
\begin{daarligbox}Han trækker ikke rigtig vejret ordentligt og ser blå ud.\end{daarligbox}
\begin{godbox}RF 28, anstrengt respiration med brug af hjælpemuskler. Central cyanose omkring læberne. SAT 88\,\% på atmosfærisk luft.\end{godbox}
\begin{hvorforbox}Beskriv RF, respirationens karakter, den objektive cyanose og den målte SAT med betingelse i stedet for ``ikke ordentligt'' og ``ser blå ud''.\end{hvorforbox}
\medskip
\begin{daarligbox}Gav ilt og det blev bedre.\end{daarligbox}
\begin{godbox}Ilt 2 l/min påbegyndt på næsekateter efter ordination; SAT stiger fra 89\,\% til 95\,\%, RF falder fra 26 til 20.\end{godbox}
\begin{hvorforbox}Angiv mængde, leveringsmåde, ``efter ordination'' og den målte effekt (SAT og RF før/efter).\end{hvorforbox}
\medskip
\begin{daarligbox}Lungerne lyder ikke så gode.\end{daarligbox}
\begin{godbox}Ved auskultation: rallen ved begge baser og pibende ekspiration bilateralt.\end{godbox}
\begin{hvorforbox}Angiv de objektive lyde og hvor de høres --- ``lyder ikke gode'' er ikke klinisk information.\end{hvorforbox}

\section{Patientnotat}
\subsection{Eksempel --- forværring, medicinsk afdeling (KOL)}
\begin{notebox}
Pt.\ 68-årig kvinde, kendt svær KOL, tiltagende åndenød gennem 2 døgn. Ved tilsyn kl.\ 14.10:
RF 28, anstrengt respiration med brug af hjælpemuskler, taler i korte sætninger. SAT 86\,\% på
atmosfærisk luft. Ved auskultation rallen ved begge baser og pibende ekspiration. Produktiv
hoste med purulent, gulligt ekspektorat. Let central cyanose. Ilt 1 l/min ordineret på
næsekateter (mål 88--92\,\%); SAT stiger til 90\,\%, RF falder til 22. Salbutamol inhalation
administreret efter ordination. Pt.\ lejret i hjertesengeleje. Lægen tilser, ordinerer i.v.\
antibiotika.
\end{notebox}

\subsection{Øvelser}
\begin{enumerate}[leftmargin=1.5em,itemsep=8pt]
\item \textbf{Omskriv objektivt:} ``Hun trak vejret helt forfærdeligt og så virkelig dårlig ud.''
\item \textbf{Omskriv:} ``Gav hende noget ilt og det hjalp lidt på iltningen.''
\item \textbf{Udfyld:} ``SAT 86\,\% \underline{\hspace{1.5cm}} atmosfærisk luft; stiger til 90\,\% \underline{\hspace{1.5cm}} 1 l ilt.''
\item \textbf{Udfyld verbet:} ``Salbutamol inhalation \underline{\hspace{2.5cm}} efter ordination.''
\item \textbf{Find vurderingen:} ``RF 28. Pt.\ er lidt en pivet type. SAT 86\,\% på atm.\ luft. Pibende ekspiration.''
\end{enumerate}

\subsection{Facit}
\begin{enumerate}[leftmargin=1.5em,itemsep=5pt]
\item RF 28, anstrengt respiration med brug af hjælpemuskler, taler i korte sætninger, let central cyanose.
\item Ilt 1 l/min ordineret på næsekateter; SAT stiger fra 86\,\% til 90\,\%.
\item \textbf{på} \ldots{} \textbf{på} (``på atmosfærisk luft'' / ``på 1 l ilt'').
\item \textbf{administreret} (sygeplejersken giver inhalationen).
\item ``Pt.\ er lidt en pivet type'' --- en nedladende, personlig vurdering.
\end{enumerate}

% ================================================================
\chapter{Smerte \& neuro}

Denne dag handler om smerte og neurologiske observationer: at beskrive smertens kvalitet,
placering og intensitet (VAS), og at vurdere bevidsthed, pupiller og orientering objektivt.
Kernen er at skelne \textbf{patientens udsagn} fra \textbf{dine observationer}.

\section{Ordforråd}
\subsection{Substantiver}
\begin{glos}
\da{smerte} & subst. & \eng{pain} & Pt.\ angiver en jagende smerte i højre side. \\
\da{VAS-skala} & subst. & \eng{VAS (pain) scale} & Smerten angives til VAS 6 ud af 10. \\
\da{udstråling} & subst. & \eng{radiation of pain} & Smerten har udstråling til venstre arm. \\
\da{smertestillende / analgetika} & subst. & \eng{analgesic} & Pt.\ får fast smertestillende og p.n.\ ved gennembrudssmerter. \\
\da{bevidsthedsniveau} & subst. & \eng{level of consciousness} & Faldende bevidsthedsniveau siden morgen. \\
\da{GCS} & subst. & \eng{Glasgow Coma Scale} & GCS 14 (E4 V4 M6). \\
\da{pupil} & subst. & \eng{pupil} & Pupiller sidelige og reagerer prompt på lys. \\
\da{orientering} & subst. & \eng{orientation} & Pt.\ er orienteret i tid, sted og person. \\
\da{hovedpine} & subst. & \eng{headache} & Pulserende hovedpine i venstre tinding. \\
\da{svimmelhed} & subst. & \eng{dizziness} & Svimmelhed ved rejsning (ortostatisk). \\
\da{kvalme} & subst. & \eng{nausea} & Kvalme ledsaget af en enkelt opkastning. \\
\da{lammelse / parese} & subst. & \eng{paralysis / paresis} & Begyndende højresidig parese i ansigt og arm. \\
\da{kramper} & subst. & \eng{seizures} & Observeret et generaliseret krampeanfald på ca.\ 40 sek. \\
\da{apopleksi} & subst. & \eng{stroke} & Pludselig lammelse rejser mistanke om apopleksi. \\
\end{glos}

\subsection{Verber}
\begin{glos}
\da{at angive (smerte)} & verbum & \eng{to report / state} & Pt.\ angiver smerte svarende til VAS 6/10. \\
\da{at lindre} & verbum & \eng{to relieve} & Smerten lindres til VAS 2/10 efter morfin. \\
\da{at vurdere (bevidsthed)} & verbum & \eng{to assess} & Bevidsthedsniveauet vurderes med GCS. \\
\da{at reagere} & verbum & \eng{to respond} & Pt.\ reagerer på tiltale, men er sløv. \\
\da{at orientere sig} & verbum & \eng{to be oriented} & Pt.\ kan ikke orientere sig i tid og sted. \\
\da{at udstråle} & verbum & \eng{to radiate} & Smerten udstråler til kæben. \\
\da{at observere (pupiller)} & verbum & \eng{to observe} & Pupilresponsen observeres hver 15.\ minut. \\
\da{at administrere} & verbum & \eng{to administer} & Morfin 5 mg administreret s.c.\ efter ordination. \\
\end{glos}

\subsection{Adjektiver}
\begin{glos}
\da{skarp / jagende} & adjektiv & \eng{sharp / stabbing} & En skarp, jagende smerte ved bevægelse. \\
\da{dump} & adjektiv & \eng{dull} & En dump smerte i lænden. \\
\da{murrende} & adjektiv & \eng{aching / nagging} & Et murrende ubehag hele dagen. \\
\da{brændende} & adjektiv & \eng{burning} & En brændende smerte langs nerven. \\
\da{pulserende} & adjektiv & \eng{throbbing} & Pulserende hovedpine, forværret af lys. \\
\da{konstant / intermitterende} & adjektiv & \eng{constant / intermittent} & Smerten er intermitterende, i anfald. \\
\da{vågen / klar} & adjektiv & \eng{awake / alert} & Pt.\ er vågen og klar ved kontakt. \\
\da{sløv / somnolent} & adjektiv & \eng{drowsy / somnolent} & Pt.\ er sløv, men vækkelig ved tiltale. \\
\da{desorienteret} & adjektiv & \eng{disoriented} & Pt.\ er desorienteret i tid og sted. \\
\da{bevidstløs} & adjektiv & \eng{unconscious} & Pt.\ er bevidstløs og reagerer ikke på smerte. \\
\end{glos}

\section{Dokumentationssprog}
\subsection{Grundregler}
\begin{enumerate}[leftmargin=1.4em,itemsep=4pt]
\item \textbf{Beskriv smerten systematisk:} placering, kvalitet (skarp/dump/brændende), intensitet (VAS), udstråling, forværrende/lindrende faktorer.
\item \textbf{Smerte er subjektiv} --- skriv ``pt.\ angiver \ldots''. Gengiv den, vurdér ikke om den er ``rigtig''.
\item \textbf{Bevidsthed} vurderes objektivt: GCS eller AVPU + orientering. Undgå ``virker lidt væk''.
\item \textbf{Pupiller} dokumenteres objektivt: størrelse, sidelighed, lysreaktion.
\item \textbf{Effekt af smertestillende} med før/efter-VAS: ``VAS 7/10 $\rightarrow$ 2/10 efter morfin 5 mg s.c.''.
\item \textbf{Ved neurologisk ændring} (parese, konfusion, kramper): beskriv objektivt, notér tidspunkt, orientér lægen.
\end{enumerate}

\subsection{Dårlig vs.\ god formulering}
\begin{daarligbox}Pt.\ har vildt ondt, klager konstant.\end{daarligbox}
\begin{godbox}Pt.\ angiver skarp, jagende smerte i højre flanke VAS 8/10 med udstråling mod lysken; forværres ved bevægelse.\end{godbox}
\begin{hvorforbox}Beskriv placering, kvalitet, intensitet (VAS), udstråling og hvad der forværrer i stedet for ``vildt ondt''.\end{hvorforbox}
\medskip
\begin{daarligbox}Han virker lidt fjern og forvirret i dag.\end{daarligbox}
\begin{godbox}Pt.\ er sløv men vækkelig ved tiltale, desorienteret i tid og sted, orienteret i person. GCS 13.\end{godbox}
\begin{hvorforbox}Brug objektive begreber: bevidsthedsniveau, orientering og GCS.\end{hvorforbox}
\medskip
\begin{daarligbox}Gav ham smertestillende, hjalp vist.\end{daarligbox}
\begin{godbox}Morfin 5 mg administreret s.c.\ kl.\ 11.20 efter ordination; VAS falder fra 7/10 til 2/10 efter 30 min.\end{godbox}
\begin{hvorforbox}Angiv præparat, dosis, vej, tid, ``efter ordination'' og den målte effekt (VAS før/efter).\end{hvorforbox}

\section{Patientnotat}
\subsection{Eksempel --- akut modtagelse (hovedpine + neuro)}
\begin{notebox}
Pt.\ 57-årig mand, pludselig kraftig hovedpine debuteret kl.\ 19.40 (``den værste nogensinde'').
Ved tilsyn kl.\ 19.55: angiver pulserende hovedpine VAS 9/10 i nakke og baghoved, ledsaget af
kvalme og en enkelt opkastning. Lyssky. Pt.\ vågen, orienteret i tid, sted og person, GCS 15.
Pupiller sidelige, reagerer prompt på lys. Ingen parese, griber lige kraftigt med begge hænder.
BT 168/96, puls 92. Paracetamol 1 g administreret p.o.\ efter ordination; VAS uændret 9/10 efter
30 min. Lægen tilser akut, ordinerer CT-skanning.
\end{notebox}

\subsection{Øvelser}
\begin{enumerate}[leftmargin=1.5em,itemsep=8pt]
\item \textbf{Omskriv objektivt:} ``Han havde helt vildt ondt i hovedet og var helt fjern.''
\item \textbf{Omskriv:} ``Gav ham en hovedpinepille men den gjorde ikke noget.''
\item \textbf{Udfyld verbet:} ``Bevidsthedsniveauet \underline{\hspace{2.5cm}} med GCS til 15.''
\item \textbf{Udfyld:} ``Pupiller sidelige og \underline{\hspace{2.5cm}} prompt på lys.''
\item \textbf{Find vurderingen:} ``VAS 9/10. Pt.\ overdriver nok lidt. Pupiller sidelige. GCS 15.''
\end{enumerate}

\subsection{Facit}
\begin{enumerate}[leftmargin=1.5em,itemsep=5pt]
\item Pt.\ angiver pulserende hovedpine VAS 9/10 i nakke/baghoved. Pt.\ vågen, orienteret i tid, sted og person, GCS 15.
\item Paracetamol 1 g administreret p.o.\ efter ordination; VAS uændret 9/10 efter 30 min.
\item \textbf{vurderes}.
\item \textbf{reagerer}.
\item ``Pt.\ overdriver nok lidt'' --- en personlig vurdering af patientens smerte.
\end{enumerate}

% ================================================================
\chapter{Medicin \& ordination}

Denne dag samler medicin-sproget: forskellen på at \textbf{ordinere} (lægen) og at
\textbf{administrere} (sygeplejersken), administrationsveje, dosis, fast vs.\ p.n.,
bivirkninger og korrekt dokumentation af given medicin.

\section{Ordforråd}
\subsection{Substantiver}
\begin{glos}
\da{ordination} & subst. & \eng{order / prescription} & Medicinen gives efter lægens ordination. \\
\da{præparat / lægemiddel} & subst. & \eng{preparation / drug} & Kontrollér, at det er det rigtige præparat. \\
\da{dosis} & subst. & \eng{dose} & Dosis er 500 mg $\times$ 3 dagligt. \\
\da{administrationsvej} & subst. & \eng{route of administration} & Administrationsvejen er peroral (p.o.). \\
\da{tablet / kapsel} & subst. & \eng{tablet / capsule} & Pt.\ indtager to tabletter med et glas vand. \\
\da{injektion} & subst. & \eng{injection} & Insulin gives som subkutan injektion. \\
\da{infusion} & subst. & \eng{infusion / drip} & Antibiotika gives som i.v.\ infusion over 30 min. \\
\da{bivirkning} & subst. & \eng{side effect} & Pt.\ observeres for bivirkninger som kvalme og svimmelhed. \\
\da{kontraindikation} & subst. & \eng{contraindication} & Kontraindikation pga.\ kendt penicillinallergi. \\
\da{p.n.-medicin} & subst. & \eng{PRN / as-needed} & Paracetamol ordineret p.n.\ ved smerter, maks.\ $\times$ 4. \\
\da{fast medicin} & subst. & \eng{regular medication} & Den faste medicin gives til morgen kl.\ 08. \\
\da{medicinliste (FMK)} & subst. & \eng{medication list} & Ordinationen fremgår af medicinlisten. \\
\da{allergi} & subst. & \eng{allergy} & Kendt allergi over for penicillin. \\
\da{gennembrudssmerte} & subst. & \eng{breakthrough pain} & P.n.-morfin gives ved gennembrudssmerte. \\
\end{glos}

\subsection{Verber}
\begin{glos}
\da{at ordinere} & verbum & \eng{to prescribe / order} & Lægen ordinerer furosemid 40 mg i.v. \\
\da{at administrere} & verbum & \eng{to administer} & Sygeplejersken administrerer medicinen efter ordination. \\
\da{at dispensere} & verbum & \eng{to dispense} & Medicinen dispenseres til doseringsæske dagen før. \\
\da{at seponere} & verbum & \eng{to discontinue} & Præparatet seponeres pga.\ bivirkninger. \\
\da{at pausere} & verbum & \eng{to withhold / pause} & Blodtryksmedicinen pauseres ved lavt BT. \\
\da{at optrappe / nedtrappe} & verbum & \eng{to titrate up / down} & Dosis nedtrappes gradvist over en uge. \\
\da{at indtage} & verbum & \eng{to take (orally)} & Pt.\ indtager tabletten selv. \\
\da{at injicere} & verbum & \eng{to inject} & Insulinen injiceres subkutant i maveskindet. \\
\da{at dobbeltkontrollere} & verbum & \eng{to double-check} & Risikomedicin dobbeltkontrolleres af to sygeplejersker. \\
\da{at observere} & verbum & \eng{to monitor} & Pt.\ observeres for bivirkninger efter første dosis. \\
\end{glos}

\subsection{Adjektiver / administrationsveje}
\begin{glos}
\da{peroral (p.o.)} & adjektiv & \eng{oral} & Tabletten gives peroralt. \\
\da{intravenøs (i.v.)} & adjektiv & \eng{intravenous} & Antibiotika gives intravenøst. \\
\da{intramuskulær (i.m.)} & adjektiv & \eng{intramuscular} & Vaccinen gives intramuskulært i overarmen. \\
\da{subkutan (s.c.)} & adjektiv & \eng{subcutaneous} & Morfin 5 mg gives subkutant. \\
\da{sublingual} & adjektiv & \eng{sublingual} & Nitroglycerin gives sublingualt. \\
\da{transdermal} & adjektiv & \eng{transdermal} & Smertestillende gives som transdermalt plaster. \\
\da{fast / p.n.} & adjektiv & \eng{regular / as-needed} & Skeln mellem fast medicin og p.n.-medicin. \\
\da{receptpligtig} & adjektiv & \eng{prescription-only} & Præparatet er receptpligtigt. \\
\end{glos}

\section{Dokumentationssprog}
\subsection{Grundregler}
\begin{enumerate}[leftmargin=1.4em,itemsep=4pt]
\item \textbf{Skeln skarpt:} LÆGEN ordinerer, SYGEPLEJERSKEN administrerer. Aldrig ``gav bare medicinen''.
\item \textbf{Dokumentér altid:} præparat, dosis, administrationsvej, tidspunkt og ``efter ordination''.
\item \textbf{Følg ``de rigtige'':} rigtige patient, præparat, dosis, tid, vej --- + dokumentation og virkning.
\item \textbf{Dokumentér virkning OG bivirkning:} ``VAS 7$\rightarrow$2 efter 30 min; ingen kvalme observeret''.
\item \textbf{Ved p.n.:} notér indikationen (``p.n.\ ved smerte VAS $\geq$ 4'') og at maksimum ikke overskrides.
\item \textbf{Hvis medicin IKKE gives:} dokumentér hvorfor (pauseret/afvist/faste) --- et tomt felt er ikke dokumentation.
\end{enumerate}

\subsection{Dårlig vs.\ god formulering}
\begin{daarligbox}Gav ham hans piller som altid.\end{daarligbox}
\begin{godbox}Fast morgenmedicin administreret p.o.\ kl.\ 08.00 efter ordination: metoprolol 50 mg, atorvastatin 40 mg.\end{godbox}
\begin{hvorforbox}Angiv præparat(er), dosis, vej, tid og ``efter ordination'' --- ``pillerne som altid'' er ikke dokumentation.\end{hvorforbox}
\medskip
\begin{daarligbox}Smertestillende p.n.\ --- gav noget morfin.\end{daarligbox}
\begin{godbox}Morfin 5 mg administreret s.c.\ p.n.\ ved gennembrudssmerte (VAS 8/10) kl.\ 13.40 efter ordination; VAS 3/10 efter 30 min., inden for maks.\end{godbox}
\begin{hvorforbox}Ved p.n.\ skal indikation, dosis, vej, tid, effekt og overholdt maksimum fremgå.\end{hvorforbox}
\medskip
\begin{daarligbox}Han ville ikke tage sin medicin.\end{daarligbox}
\begin{godbox}Pt.\ afviser aftenmedicin kl.\ 20; angiver kvalme. Informeret om formålet. Medicin ikke givet, lægen orienteret.\end{godbox}
\begin{hvorforbox}Ved afvisning: dokumentér handlingen, patientens begrundelse i egne ord, din information og at lægen er orienteret.\end{hvorforbox}

\section{Patientnotat}
\subsection{Eksempel --- medicinadministration, medicinsk afdeling}
\begin{notebox}
Medicingivning morgen kl.\ 08.00: Fast medicin administreret p.o.\ efter ordination --- metoprolol
50 mg, atorvastatin 40 mg, pantoprazol 40 mg. Identitet kontrolleret mod medicinlisten. BT før
metoprolol 138/84, puls 76, derfor givet. Insulin (Novorapid 8 IE) administreret subkutant efter
ordination, dobbeltkontrolleret af kollega. Kl.\ 13.40: Pt.\ angiver gennembrudssmerte VAS 8/10;
morfin 5 mg administreret s.c.\ p.n.\ efter ordination; VAS 3/10 efter 30 min., ingen
respirationspåvirkning. Amlodipin pauseret til aften pga.\ BT 96/58 efter aftale med lægen;
genoptages ved systolisk BT $>$ 110. Pt.\ observeret for bivirkninger, ingen bemærket.
\end{notebox}

\subsection{Øvelser}
\begin{enumerate}[leftmargin=1.5em,itemsep=8pt]
\item \textbf{Omskriv fagligt:} ``Gav ham hans sædvanlige morgenpiller.''
\item \textbf{Omskriv:} ``Han havde ondt, så jeg gav lidt morfin, det hjalp.''
\item \textbf{Udfyld verberne:} ``Lægen \underline{\hspace{2cm}} furosemid 40 mg i.v.; sygeplejersken \underline{\hspace{2.5cm}} det efter ordination.''
\item \textbf{Udfyld:} ``Amlodipin \underline{\hspace{2cm}} til aften pga.\ BT 96/58 efter aftale med lægen.''
\item \textbf{Find fejlen:} ``Gav noget vanddrivende. Virkede vist. Pt.\ lidt besværlig.''
\end{enumerate}

\subsection{Facit}
\begin{enumerate}[leftmargin=1.5em,itemsep=5pt]
\item Fast morgenmedicin administreret p.o.\ efter ordination kl.\ 08.00: metoprolol 50 mg, atorvastatin 40 mg, pantoprazol 40 mg.
\item Pt.\ angiver gennembrudssmerte VAS 8/10; morfin 5 mg administreret s.c.\ p.n.\ efter ordination; VAS 3/10 efter 30 min.
\item \textbf{ordinerer} \ldots{} \textbf{administrerer}.
\item \textbf{pauseret}.
\item Mangler præparat/dosis/vej/tid/``efter ordination'' og målt effekt; desuden er ``lidt besværlig'' en personlig vurdering.
\end{enumerate}

\vfill
\noindent{\small\color{muted}\rule{\linewidth}{0.4pt}\\[4pt]
\textbf{Sådan tilføjes en ny dag:} kopiér et kapitel, ret systemet, og genbrug de fire sektioner
(ordforråd, dokumentation, patientnotat, øvelser). Web-appens datafil
\texttt{assets/data/dagN\_\ldots.js} og kapitlet holdes med samme indhold.}

\end{document}
